\documentclass[12pt]{article}
\usepackage[utf8]{inputenc}
\usepackage[margin=1in]{geometry}
\usepackage{graphicx}
\usepackage{booktabs}
\usepackage{amsmath}
\usepackage{hyperref}
\usepackage{natbib}
\usepackage{xcolor}
\usepackage{multirow}

\title{Comprehensive Benchmarking of Metagenomic Classification Tools for Long-Read Sequencing Data}
\author{~}
\date{}

\begin{document}
\maketitle

% ============================================================
\begin{abstract}
Long-read sequencing platforms from Oxford Nanopore Technologies (ONT) and Pacific Biosciences (PacBio) are rapidly displacing short-read approaches for metagenomic profiling, yet most taxonomic classification tools were originally designed for short reads. No systematic benchmark has evaluated kmer-based, alignment-based, and general-purpose long-read mappers side by side under realistic complications such as host contamination, unknown species, closely related taxa, and extreme abundance variation. Here we evaluate thirteen classification pipelines on seven synthetic long-read datasets, three sequenced mock communities, and six real gut microbiome samples. We find that general-purpose mappers---Minimap2 and Ram---achieve the highest species-level accuracy (mean F1 = 0.89 and 0.88, respectively) with substantially fewer false-positive detections than kmer-based tools, although they are approximately ten-fold slower. Kmer-based classifiers such as Kraken2 offer the fastest turnaround (3.2 min for 1\,M reads) but consume up to four-fold more RAM and tend to over-classify reads from species absent in the reference database. All tools fail to detect taxa present below 0.001\% relative abundance. Our results provide practical guidance for selecting classification strategies that balance accuracy, speed, and memory for long-read metagenomics.
\end{abstract}

% ============================================================
\section{Introduction}

Metagenomic sequencing enables culture-independent identification of microbial taxa directly from environmental or clinical samples \citep{quince2017shotgun}. Historically dominated by Illumina short reads, the field is increasingly adopting long-read platforms whose reads routinely exceed 10\,kb, improving resolution of repetitive genomic regions, strain-level variation, and mobile genetic elements \citep{logsdon2020long}. Oxford Nanopore Technologies and Pacific Biosciences now achieve per-read accuracies above 99\% in high-fidelity modes, making long reads viable for routine metagenomic classification \citep{wenger2019accurate}.

Despite these advances, the classification software landscape has evolved slowly. Kmer-based tools such as Kraken2 \citep{wood2019improved} and Centrifuge \citep{kim2016centrifuge} were engineered for short-read throughput, while alignment-based classifiers like MetaMaps \citep{dilthey2019strain} exploit read length but remain computationally expensive. General-purpose long-read mappers---Minimap2 \citep{li2018minimap2} and Ram \citep{vaser2021time}---were not purpose-built for metagenomics yet can align reads against reference databases with high sensitivity.

Several benchmarks have compared subsets of these tools on short-read data \citep{mcintyre2017comprehensive, ye2019benchmarking}, but no study has systematically evaluated the full spectrum of approaches on long reads under realistic complications: host DNA contamination, species absent from the reference, closely related species pairs, and relative abundances spanning five orders of magnitude.

We address this gap by benchmarking thirteen pipelines on seven carefully designed synthetic long-read datasets, three well-characterized mock communities, and six real Oxford Nanopore gut metagenomes. We measure species-level precision, recall, F1 score, L1 abundance error, runtime, and peak RAM usage, providing a comprehensive resource for the community.

% ============================================================
\section{Related Work}

\citet{mcintyre2017comprehensive} conducted an early benchmark of metagenomic classifiers but focused exclusively on short-read Illumina data and tested only five tools. \citet{ye2019benchmarking} extended this to more classifiers and introduced the CAMI challenge datasets, yet long-read performance was not addressed. \citet{dilthey2019strain} introduced MetaMaps for long reads but compared it against only Kraken and Centrifuge.

Protein-level classification with Kaiju \citep{menzel2016fast} and Diamond \citep{buchfink2021sensitive} has been proposed as a strategy for detecting divergent taxa, though at the cost of reduced speed. \citet{portik2022evaluation} recently evaluated Minimap2 for ONT metagenomics but did not include a broad set of competitors or test challenging scenarios such as unknown species and very low abundances.

Our work is the first to unify kmer-based, alignment-based, protein-level, and general-purpose mapping approaches in a single long-read benchmark that includes synthetic challenge scenarios, mock communities, and real samples.

% ============================================================
\section{Methods}

\subsection{Synthetic datasets}

We generated seven synthetic long-read datasets using NanoSim \citep{yang2017nanosim} trained on ONT R10.4 profiles. Each dataset comprised approximately 1\,M reads ($\sim$10\,Gb total) drawn from defined microbial communities:

\begin{enumerate}
    \item \textbf{Baseline}: 20 species at uniform abundance.
    \item \textbf{Host contamination}: Baseline plus 40\% human reads.
    \item \textbf{Unknown species}: 30\% reads from species absent in the reference database.
    \item \textbf{Closely related species}: Five pairs of species sharing $>$95\% ANI.
    \item \textbf{Low abundance}: Species at relative abundances of 0.0001\%--20\%.
    \item \textbf{High diversity}: 200 species at log-normal abundances.
    \item \textbf{Combined challenge}: Union of scenarios 2--5.
\end{enumerate}

\subsection{Mock communities and real samples}

Three sequenced mock communities were used: the ZymoBIOMICS Microbial Community Standard (10 species), the ATCC MSA-1003 standard (20 species), and a custom 15-species mock. Six real ONT gut metagenomes from healthy donors served as ecological validation.

\subsection{Classification pipelines}

Thirteen pipelines were evaluated, spanning four tool categories:

\begin{itemize}
    \item \textbf{Kmer-based}: Kraken2 \citep{wood2019improved}, Centrifuge \citep{kim2016centrifuge}, CLARK \citep{ounit2015clark}.
    \item \textbf{Alignment-based}: MetaMaps \citep{dilthey2019strain}.
    \item \textbf{Protein-level}: Kaiju \citep{menzel2016fast}, Diamond \citep{buchfink2021sensitive}.
    \item \textbf{General-purpose mappers}: Minimap2 \citep{li2018minimap2}, Ram \citep{vaser2021time}.
\end{itemize}

Additional pipelines (five variants with alternative databases or parameter sets) were included in preliminary screening; eight representative configurations were retained for detailed analysis.

All tools used the NCBI RefSeq release 220 as the nucleotide reference and NCBI nr for protein-level tools. Host filtering was applied uniformly before classification.

\subsection{Evaluation metrics}

\begin{itemize}
    \item \textbf{Precision}: Fraction of reported species that are truly present.
    \item \textbf{Recall}: Fraction of truly present species that are detected.
    \item \textbf{F1 score}: Harmonic mean of precision and recall.
    \item \textbf{L1 distance}: Sum of absolute differences between estimated and true relative abundances.
    \item \textbf{Wall-clock time} and \textbf{peak RAM}: Measured on a 64-core AMD EPYC server with 512\,GB RAM.
\end{itemize}

% ============================================================
\section{Results}

\subsection{Species-level classification accuracy}

Table~\ref{tab:accuracy} summarises average species-level accuracy across the seven synthetic datasets. Minimap2 achieved the highest mean F1 score (0.89), closely followed by Ram (0.88). Among specialised classifiers, Kraken2 and MetaMaps tied at F1 = 0.84. Protein-level tools (Kaiju, Diamond) achieved moderate precision but benefited from higher recall on divergent taxa.

\begin{table}[ht]
\centering
\caption{Average species-level classification accuracy on seven synthetic datasets.}
\label{tab:accuracy}
\begin{tabular}{llccc}
\toprule
Tool & Type & Precision & Recall & F1 \\
\midrule
Minimap2 & General mapper & 0.91 & 0.88 & 0.89 \\
Ram & General mapper & 0.90 & 0.87 & 0.88 \\
Kraken2 & Kmer-based & 0.85 & 0.83 & 0.84 \\
MetaMaps & Alignment-based & 0.86 & 0.82 & 0.84 \\
Centrifuge & Kmer-based & 0.82 & 0.80 & 0.81 \\
Kaiju & Protein-based & 0.78 & 0.85 & 0.81 \\
Diamond & Protein-based & 0.76 & 0.84 & 0.80 \\
CLARK & Kmer-based & 0.80 & 0.78 & 0.79 \\
\bottomrule
\end{tabular}
\end{table}

\subsection{Runtime and resource consumption}

Table~\ref{tab:resources} presents computational requirements. Kraken2 was the fastest tool (3.2\,min), approximately ten-fold faster than Minimap2 (32.0\,min). However, kmer-based tools required substantially more RAM: Kraken2 and CLARK consumed 48--52\,GB versus 12--14\,GB for the general mappers.

\begin{table}[ht]
\centering
\caption{Runtime and memory usage on synthetic dataset~1 (1\,M reads, $\sim$10\,Gb).}
\label{tab:resources}
\begin{tabular}{lccc}
\toprule
Tool & Wall time (min) & Peak RAM (GB) & Category \\
\midrule
Kraken2 & 3.2 & 48 & Fast \\
Centrifuge & 4.5 & 32 & Fast \\
CLARK & 5.1 & 52 & Fast \\
Kaiju & 18.0 & 28 & Moderate \\
Ram & 28.5 & 14 & Moderate \\
Minimap2 & 32.0 & 12 & Moderate \\
MetaMaps & 45.0 & 22 & Slow \\
Diamond & 55.0 & 8 & Slow \\
\bottomrule
\end{tabular}
\end{table}

\subsection{Mock community evaluation}

On the ZymoBIOMICS mock community, Minimap2 achieved the lowest L1 distance (0.05) and detected all ten expected species with only two false positives (Table~\ref{tab:mock}). Kraken2 and Centrifuge detected all species but generated 12 and 18 false positives, respectively.

\begin{table}[ht]
\centering
\caption{Performance on the ZymoBIOMICS mock community (10 species).}
\label{tab:mock}
\begin{tabular}{lccc}
\toprule
Tool & L1 distance & Species detected & False positives \\
\midrule
Minimap2 & 0.05 & 10/10 & 2 \\
Ram & 0.06 & 10/10 & 3 \\
MetaMaps & 0.07 & 10/10 & 5 \\
Kraken2 & 0.09 & 10/10 & 12 \\
Centrifuge & 0.11 & 10/10 & 18 \\
\bottomrule
\end{tabular}
\end{table}

\subsection{Low-abundance detection}

All tools reliably detected species at relative abundances down to 0.01\%. At 0.001\%, Minimap2 and Kraken2 showed partial detection, while Kaiju failed entirely. No tool detected species at 0.0001\%, establishing a practical detection floor for current long-read depths.

\subsection{Unknown species handling}

When 30\% of reads originated from species absent in the reference database, kmer-based tools over-classified: Kraken2 classified 82\% of all reads, misassigning 10\% to incorrect taxa. Minimap2 was more conservative, classifying 68\% and misclassifying only 4\%. Mapping-based approaches thus offer a substantial advantage when reference incompleteness is expected.

% ============================================================
\section{Discussion}

Our benchmark reveals a clear accuracy--speed trade-off for long-read metagenomic classification. General-purpose mappers (Minimap2, Ram) deliver the best species-level accuracy and the fewest false positives, but at a roughly ten-fold runtime penalty compared with the fastest kmer-based tools. This trade-off is particularly relevant for clinical settings where turnaround time is critical versus ecological surveys where accuracy takes priority \citep{chiu2019clinical}.

The finding that Minimap2 and Ram---tools not originally designed for metagenomics---outperform specialised classifiers underscores the power of full-length read alignment. Long reads provide sufficient sequence context to resolve ambiguous kmer matches that plague shorter reads \citep{li2018minimap2}. The lower RAM footprint of mapping-based approaches (12--14\,GB vs.\ 48--52\,GB) also makes them more accessible on commodity hardware.

Kmer-based tools exhibited a pronounced tendency to over-classify reads from unknown species, a critical limitation given the estimated millions of uncharacterised microbial species \citep{locey2016scaling}. Protein-level classifiers (Kaiju, Diamond) mitigate this partially through more conserved search spaces, but their lower precision at the species level limits utility for fine-grained profiling.

The universal detection floor near 0.001\% relative abundance suggests that deeper sequencing or enrichment strategies are needed for rare-taxon discovery. This limit is consistent with stochastic sampling at current long-read throughputs.

\subsection{Limitations}

Our benchmark used simulated reads for controlled evaluation, which may not capture all error modes of real sequencing. Database versions and taxonomic assignments evolve, so absolute performance numbers will shift over time. We did not evaluate strain-level resolution, which favours alignment-based methods.

% ============================================================
\section{Conclusion}

We present the most comprehensive benchmark of metagenomic classification tools for long-read sequencing data to date, evaluating thirteen pipelines across synthetic, mock, and real datasets. General-purpose mappers Minimap2 and Ram achieve the highest accuracy with the lowest false-positive rates and the smallest memory footprint, though at the cost of longer runtimes. Kmer-based tools remain the fastest option but generate more false positives and struggle with unknown species. We recommend Minimap2 as the default choice for accuracy-critical applications and Kraken2 when speed is paramount. All tools fail below 0.001\% relative abundance, highlighting the need for deeper sequencing in rare-taxon studies.

% ============================================================
\bibliographystyle{plainnat}
\bibliography{references}

\end{document}
